\section{Appendix D. Negative Results and Hardened Engineering}

One of the primary contributions of HIBA is its production-hardened engineering, born directly from extensive failure analysis. We document these negative results to aid future researchers.

\subsection{D.1 The RoBERTa Collapse (Nested Encoder Attenuation)}

\textbf{Symptom:} When injecting early versions of block-sparse attention into RoBERTa, the model accuracy catastrophically collapsed to $\approx 50.92\%$ (random chance on SST-2), while DistilBERT maintained $>90\%$.

\textbf{Investigation:} 
We traced the gradient flow and discovered a fundamental structural difference in HuggingFace's implementation of the BERT family. DistilBERT uses a "flat" architecture where \texttt{out\_lin} is inside the \texttt{MultiHeadSelfAttention} block. 
RoBERTa uses a "nested" architecture:
\begin{lstlisting}
class RobertaAttention:
    self: RobertaSelfAttention      # Contains Q, K, V
    output: RobertaSelfOutput       # Contains Dense (Output Projection)

\end{lstlisting}

\textbf{Why it failed:}
Our initial module applied its own \texttt{out\_proj} and returned the result to \texttt{RobertaAttention}. RoBERTa then passed this output into \texttt{RobertaSelfOutput.dense}. This resulted in the signal passing through two consecutive dense layers without a non-linearity: \texttt{Dense(Dense(x))}. This double-projection drastically altered the variance of the residual stream, causing immediate gradient explosion/vanishing and representation collapse.

\textbf{The Fix:}
We introduced explicit introspection (\texttt{self.is\_nested\_encoder}) during \texttt{copy\_weights}. If a nested architecture is detected, HIBA extracts the inner $Q, K, V$ weights, but \textbf{bypasses its own \texttt{out\_proj} during the forward pass}, returning the raw context matrix to be safely processed by RoBERTa's native outer wrapper. This completely resolved the collapse, bringing RoBERTa's accuracy to 78.40\% ($\Delta$ = -0.16).

\subsection{D.2 The Qwen/LLaMA RoPE Mismatch}

\textbf{Symptom:} Injecting HIBA into Qwen3.5 resulted in NaNs and severe language modeling degradation.

\textbf{Investigation:}
Standard implementations of Rotary Position Embeddings (RoPE) assume that the entire head dimension ($d_h$) is rotated. However, modern models often apply partial rotations to save compute. For instance, Qwen3.5 uses $d_h = 256$ but only rotates the first \texttt{rotary\_dim = 64} coordinates.

\textbf{The Fix:}
We rewrote the \texttt{apply\_rotary\_pos\_emb} function to dynamically slice the query and key tensors:

\begin{lstlisting}
if rotary_dim < head_dim:
    q_rot, q_pass = q[..., :rotary_dim], q[..., rotary_dim:]
    # Rotate q_rot, then torch.cat with q_pass

\end{lstlisting}

This partial-dimension splicing restored numerical parity with the base Qwen3.5 forward pass up to $1e^{-5}$ tolerance.

\subsection{D.3 Failed Attempts at Pure Sparsity}

Before implementing \textbf{Periodic Dense Rectification}, we attempted to train HIBA as a purely sparse model (\texttt{rectification\_interval = infinity}). While loss decreased during training, the model failed to generalize on validation data involving complex negations (e.g., "Not exactly a very good movie"). The strict block boundary severed the connection between "Not" and "good". The introduction of the dense rectification layer acts as an "attentional highway," allowing these critical long-range linguistic modifiers to route properly.


